\section{Network Model Validation}

\subsection{Signal Level Validation}

To establish that the simulated ns3 channel reproduces physically meaningful propagation rather than arbitrary values, the received signal power reported by ns3 was validated against the Friis free space path propagation loss model. Received power was obtained at five UAV-ground station separations spanning $58.5 - 113.8m$, comprising $150,393$ packet samples and the per position mean was compared against the predictions obtained at the same distance using the Friis propagation model. The two agree closely. Across the five positions the comparison yields a coefficient of determination of $R^2 = 0.86$, a root mean squared error of $0.77dB$ and a mean bias of $-0.77dB$. The measured points follow the 1:1 agreement line, confirming that the simulated signal reproduces the expected distance dependent decay. No trend could be identified with the residuals. The error is an almost constant offset of about $-0.8dB$ rather than a distance dependent deviation, indicating that ns3 reports received power roughly $0.77dB$ below the analytical value uniformly across the tested range. The channel model therfore reproduces the shape of free space propagation with a slight, fixed offset in absolute power that could be calibrated out where exact levels are required. This check covered a single five position session and is treated as a trend-level validation of the propagation model rather than a set of independent statistical replications.


\subsection{Obstacle Aware Attenuation}

Without limiting to free space path loss, the channel model adds a per link attenuation term whenever the line of sight between two UAV is obstructed by world geometry. To confirm this mechanism behaves as expected, the per link obstacle loss was recorded throughout a multi UAV mission run and examined against the UAVs' positions relative to the buildings in the physical environment. The result matches the expected geometry. Almost every link incurs obstacle loss at some point during the mission as the UAVs move behind buildings and the loss appears and clears in step with line of sight rather than at random. The \emph{UAV1 - UAV2} link shows the highest attenuation which the reason is solid since these two vehicles operates on opposite sides of several buildings for part of the mission. This confirms that the obstacle term responds to the actual physical environment giving the simulated channel a spatially varying, environment dependent quality that a distance only model would ideally miss. This property lets the framework capture how the physical environment changes considering not only separation, shapes multi UAV communication. This check is descriptive, it demonstrates that the obstacle mechanism engages correctly with scene geometry and is not offered as a quantitative accuracy claim.